\section{Limitations and failure criteria}
\label{sec:limitations}

\subsection{Limitations}

\paragraph{No calibration.} Not one coefficient in
Eqs.~\eqref{eq:dI}--\eqref{eq:dO} is constrained by data, because no relevant
data exist. The model is dimensionless and exploratory. It supports statements
of the form ``if these mechanisms operate with these relative strengths, then
this follows'', and it supports nothing else. The numbers in
\S\ref{sec:results} illustrate qualitative regimes. They are not estimates of
anything, and should never be quoted as though they were.

\paragraph{Collapse is a proxy, not a trajectory.} There is no mortality term.
Stage~4a therefore raises the collapse criteria --- regulation below $\Rmin$,
amplification exceeding regulation by more than $\Theta$ --- while capability
carries on quite happily. The ``transient spike, then nothing'' signature people
associate with civilizational collapse is consequently \emph{not} reproduced
dynamically. Coupling mortality to regulatory exhaustion is the obvious fix and
we have not done it.

\paragraph{Stiffness.} Making suppression fractional costs something numerically.
The relaxation rate $\Sigma$ scales as $\Icap^{2}\Reg$ under the superlinear
variants, and an explicit method needs step sizes of order $1/\Sigma$, so
integration cost grows linearly with the suppression weights. That bounds how
far into the strongly-suppressing corner of parameter space we can economically
go. An implicit or stiffness-switching integrator would remove the limit.

\paragraph{The Lambert boundary is auxiliary.} As \S\ref{sec:lambert} establishes,
$C_R$ does not appear in the integrated system. Nothing in this paper depends on
the Lambert boundary, and it must not be read as a bifurcation of the dynamics.

\paragraph{Quietness is degenerate.} Deliberate stealth and emergent efficiency
enter $\Sigma$ in exactly the same way, so observability alone cannot tell them
apart. A civilization hiding on purpose and one that simply got very good at
compression look identical from here.

\paragraph{No feedback from capability to regulation.} Equation~\eqref{eq:dR}
makes regulatory capacity depend on the exogenous drivers and on erosion by
$\Arec$, but not on $\Icap$ at all (Figure~\ref{fig:schematic}). This is a real
simplification and probably the most consequential one. A highly capable
civilization would plausibly be better at building governance as well as better
at eroding it, and which of those wins is not obvious. Closing the loop --- say
by making the building term depend on $\Aref \Icap$ --- is straightforward, and
it would open up genuine bistability between well-regulated and poorly-regulated
high-capability states. We think that is the most interesting thing anyone could
do with this model next.

\paragraph{One civilization, not a population.} The model integrates a single
trajectory. Any statement about $\Nobs$ therefore assumes a population whose
members all share parameters, which is plainly false. A population-level
treatment would integrate over distributions of the drivers and the branch
probabilities, and it is required before this framework can be compared
quantitatively with grabby-aliens arguments
\citep{hanson2021grabby,prantzos2020}.

\subsection{Failure criteria}

A model of this kind is hard to argue with unless its author says in advance
what would count against it. So:

\begin{enumerate}\itemsep3pt
  \item \textbf{Non-robustness.} If the low-observability regime arises only for
    a narrow band of suppression scalings, or only under one choice from
    Table~\ref{tab:registry}, then the conclusion is an artefact of that choice.
    \S\ref{sec:results} reports the dependence explicitly; the regime does not
    arise under observability model~A.
  \item \textbf{Assumed rather than derived decline.} If peak-then-decline could
    be obtained only by selecting a peaked $\Obs(\Icap)$, this framework would
    add nothing to the low-observability arguments that already exist.
    Equation~\eqref{eq:decline} and \S\ref{sec:derived} are our answer. Show
    that the required scaling separation --- $\Sigma$ growing faster in $\Icap$
    than $\Pi$ --- is physically implausible, and the answer fails.
  \item \textbf{Collapse and quietness not separated.} If the model could not
    distinguish a civilization that ended from one that went quiet, it would
    reproduce the exact ambiguity it was built to dissolve. The regime criteria
    do separate them, but only as far as the collapse proxy means anything --- see
    above.
  \item \textbf{Thermodynamic violation.} Any result requiring $\Obs \leq 0$
    invalidates the physical interpretation outright.
    Proposition~\ref{prop:obs} makes that unreachable by construction, and the
    property-based tests of \S\ref{sec:verification} check it empirically over
    300 randomised configurations.
  \item \textbf{Search coverage ignored.} Drop $\Psearch$ and non-detection gets
    misread as evidence about $\Ntrue$. It appears explicitly in
    Eq.~\eqref{eq:decomp} and must stay there.
  \item \textbf{Empirical pressure.} Sustained null results from infrared and
    waste-heat surveys at improving sensitivity would argue against the quiet
    regime being populated, since that is precisely the channel it predicts
    stays open \citep{griffith2015,garrett2015,wright2022haystack}.
\end{enumerate}

Only the last of those is empirical rather than internal. That asymmetry is a
limitation of the subject rather than of the model, and we would rather name it
than dress it up.
